\part{Per Iniziare}

\chapter{Introduzione}

FairWindSK è un sistema di visualizzazione nautica per plancia costruito attorno a Signal~K.
Riunisce dati nautici in tempo reale, applicazioni di bordo e strumenti operativi nativi
in un'unica interfaccia stabile progettata per il marinaio al timone.

\begin{notebox}
FairWindSK fa parte del programma di ricerca DYNAMO dell'Università degli Studi di Napoli
'Parthenope', ora sostenuto dal progetto DataX4Sea finanziato da NEC Laboratories America.
L'obiettivo è il crowdsourcing di dati ambientali costieri tramite flotte di imbarcazioni da diporto.
Tutti i risultati sono open-source e open-data.
\end{notebox}

\section{Che cosa fa FairWindSK}

FairWindSK è una \emph{shell} di visualizzazione, non un server di dati.
Avvolge un server Signal~K in esecuzione, ne scopre le applicazioni web installate e le
presenta in un frame stabile insieme ai controlli operativi nativi.

Pensa a FairWindSK come alla cornice attorno alle app di bordo e ai dati di navigazione
in tempo reale: ti aiuta a tenere le informazioni chiave sempre visibili mentre passi
da un'attività all'altra.

\subsection{Piattaforme supportate}

\rowcolors{2}{LtGray}{white}
\begin{longtable}{L{3.0cm}L{4.8cm}L{5.9cm}}
\toprule
\rowcolor{Navy}
\color{white}\sffamily\bfseries Piattaforma &
\color{white}\sffamily\bfseries Utilizzo tipico &
\color{white}\sffamily\bfseries Note \\
\midrule
Raspberry Pi OS & Display di plancia dedicato, avvio chiosco &
  Funzionalità desktop completa; strumenti kiosk in \code{extras/} \\
Linux           & PC o laptop da postazione nautica    & Funzionalità desktop completa \\
macOS           & MacBook da postazione nautica         & Funzionalità desktop completa \\
Windows         & PC da postazione nautica              & Funzionalità desktop completa \\
Android         & Tablet di plancia, display secondario & Implementazione mobile Qt WebView \\
iOS / iPadOS    & iPad di plancia o display portatile   & Implementazione mobile Qt WebView \\
\bottomrule
\end{longtable}

\section{Cosa NON è FairWindSK}

FairWindSK dipende interamente da un server Signal~K in esecuzione sulla rete di bordo.
Senza Signal~K non ha dati da visualizzare, nessuna app da avviare e nessuno strumento
da leggere. Non elabora direttamente le sentenze NMEA, non si collega direttamente agli
strumenti e non archivia le carte nautiche.

\begin{warningbox}
FairWindSK è un ausilio alla navigazione.
Non sostituisce la perizia nautica, la guardia visiva, le carte nautiche cartacee
o le procedure di sicurezza certificate.
Il comandante dell'imbarcazione è il solo responsabile della condotta in sicurezza
della traversata in ogni momento.
\end{warningbox}

\section{Cosa serve prima di iniziare}

\begin{itemize}
  \item Un server Signal~K in esecuzione e raggiungibile sulla rete Wi-Fi di bordo.
        Il server deve esporre l'API standard su \code{<server>/signalk} e il catalogo
        app su \code{<server>/signalk/v1/apps/list}.
  \item I propri strumenti (GPS, vento, profondità, AIS, autopilota, monitor motore)
        che alimentano Signal~K via NMEA~0183, NMEA~2000 o SeaTalk.
  \item FairWindSK installato su un dispositivo di visualizzazione.
\end{itemize}

\section{Indicazioni per il dispositivo}

\subsection{Tablet e display di plancia}

I tablet e i display nautici di grandi dimensioni sono la soluzione migliore per l'uso
primario di plancia: target di tocco più grandi, layout stabile, lettura facile a distanza
di braccio e maggiore spazio per la consapevolezza simultanea di rotta e app.

\subsection{Smartphone}

Gli smartphone funzionano come display secondari. Il layout sarà più compatto e ci sarà
meno spazio per la consapevolezza simultanea. In fasi di navigazione ravvicinata, manovre
di recupero o situazioni ad alto carico di lavoro, un display più grande è più sicuro.

\subsection{Sistemi desktop}

Su piattaforme desktop FairWindSK può essere usato come display fisso da postazione
nautica. L'interfaccia va trattata come un display operativo nautico prima di tutto,
indipendentemente dal metodo di input.

\chapter{Prima di Lasciare la Banchina}

Esegui questi controlli prima di ogni partenza.
Cinque minuti in banchina valgono un'ora in mare.

\section{Configurazione iniziale}

Quando FairWindSK è installato su un nuovo display o su un sistema appena configurato,
completa questi passaggi prima di affidarti al sistema in navigazione:

\begin{enumerate}
  \item Apri \ui{Impostazioni~> Connessione}.
        Inserisci l'indirizzo del server Signal~K, ad es.\ \code{http://192.168.1.100:3000}.
        Tocca \key{Connetti}.
        Attendi che lo stato diventi \textit{Live}.
  \item Apri \ui{Impostazioni~> Unità}.
        Scegli nodi per la velocità, miglia nautiche per la distanza, metri o piedi per la profondità.
  \item Apri \ui{Impostazioni~> Applicazioni}.
        Verifica che le app che ti aspetti siano presenti.
        Organizza la pagina~1 del launcher con le app che usi in ogni traversata.
  \item Apri \ui{Impostazioni~> Comfort} e scegli il preset \textit{Giorno}.
        Verifica la leggibilità dalla posizione di plancia.
  \item Se l'installazione include autopilota, ancora, allarmi o altre funzioni
        basate su plugin, conferma che ciascuna sia disponibile e risponda correttamente
        prima della partenza.
  \item Se il dispositivo non ha tastiera fisica, abilita la tastiera virtuale in
        \ui{Impostazioni~> Principale} e riavvia FairWindSK.
  \item Esegui una breve traversata di prova in condizioni calme prima di affidarti
        a FairWindSK in mare aperto.
\end{enumerate}

\begin{warningbox}
Non dare per scontato che una funzione sia pronta solo perché la sua icona è visibile.
Conferma prima i dati in tempo reale, la risposta ai comandi e il comportamento reale.
\end{warningbox}

\section{Lista di controllo pre-partenza}

\rowcolors{2}{LtGray}{white}
\begin{longtable}{C{0.8cm}L{12.9cm}}
\toprule
\rowcolor{Navy}
\color{white}\sffamily\bfseries $\square$ &
\color{white}\sffamily\bfseries Controllo \\
\midrule
$\square$ & FairWindSK si avvia normalmente e lo stato della shell è normale \\
$\square$ & Connessione Signal~K integra (\textit{Live} nell'area di stato) \\
$\square$ & Posizione viva nella Barra Superiore --- latitudine e longitudine plausibili \\
$\square$ & SOG in aggiornamento --- il GPS ha un fix valido \\
$\square$ & Profondità in aggiornamento (se il ecoscandaglio è installato) \\
$\square$ & Rotta (HDG) visualizzata da bussola o GPS \\
$\square$ & Plotter aperto e imbarcazione nella posizione corretta sulla carta \\
$\square$ & Rotta o waypoint attivo se si procede verso una destinazione \\
$\square$ & Target AIS visibili nel plotter (se il ricevitore AIS è installato) \\
$\square$ & Luminosità dello schermo e preset comfort adeguati alle condizioni di luce \\
$\square$ & Schermo leggibile e raggiungibile dalla normale posizione operativa \\
$\square$ & Posizione del tasto POB confermata --- sapere esattamente dov'è \\
$\square$ & Autopilota disponibile e funzionante (se richiesto) \\
$\square$ & Funzione guardia ancora disponibile (se si prevede di ancorare) \\
\bottomrule
\end{longtable}
